% Options for packages loaded elsewhere
\PassOptionsToPackage{unicode}{hyperref}
\PassOptionsToPackage{hyphens}{url}
\PassOptionsToPackage{dvipsnames,svgnames,x11names}{xcolor}
\PassOptionsToPackage{space}{xeCJK}
\documentclass[
]{article}
\usepackage{xcolor}
\usepackage[margin=2.5cm]{geometry}
\usepackage{amsmath,amssymb}
\setcounter{secnumdepth}{-\maxdimen} % remove section numbering
\usepackage{iftex}
\ifPDFTeX
  \usepackage[T1]{fontenc}
  \usepackage[utf8]{inputenc}
  \usepackage{textcomp} % provide euro and other symbols
\else % if luatex or xetex
  \usepackage{unicode-math} % this also loads fontspec
  \defaultfontfeatures{Scale=MatchLowercase}
  \defaultfontfeatures[\rmfamily]{Ligatures=TeX,Scale=1}
\fi
\usepackage{lmodern}
\ifPDFTeX\else
  % xetex/luatex font selection
  \setmainfont[]{DejaVu Sans}
  \setmonofont[]{DejaVu Sans Mono}
  \ifXeTeX
    \usepackage{xeCJK}
    \setCJKmainfont[]{Noto Sans CJK SC}
  \fi
  \ifLuaTeX
    \usepackage[]{luatexja-fontspec}
    \setmainjfont[]{Noto Sans CJK SC}
  \fi
\fi
% Use upquote if available, for straight quotes in verbatim environments
\IfFileExists{upquote.sty}{\usepackage{upquote}}{}
\IfFileExists{microtype.sty}{% use microtype if available
  \usepackage[]{microtype}
  \UseMicrotypeSet[protrusion]{basicmath} % disable protrusion for tt fonts
}{}
\makeatletter
\@ifundefined{KOMAClassName}{% if non-KOMA class
  \IfFileExists{parskip.sty}{%
    \usepackage{parskip}
  }{% else
    \setlength{\parindent}{0pt}
    \setlength{\parskip}{6pt plus 2pt minus 1pt}}
}{% if KOMA class
  \KOMAoptions{parskip=half}}
\makeatother
\usepackage{longtable,booktabs,array}
\usepackage{caption}
\captionsetup[table]{skip=6pt}
\newcounter{none} % for unnumbered tables
\usepackage{calc} % for calculating minipage widths
% Correct order of tables after \paragraph or \subparagraph
\usepackage{etoolbox}
\makeatletter
\patchcmd\longtable{\par}{\if@noskipsec\mbox{}\fi\par}{}{}
\makeatother
% Allow footnotes in longtable head/foot
\IfFileExists{footnotehyper.sty}{\usepackage{footnotehyper}}{\usepackage{footnote}}
\makesavenoteenv{longtable}
\setlength{\emergencystretch}{3em} % prevent overfull lines
\providecommand{\tightlist}{%
  \setlength{\itemsep}{0pt}\setlength{\parskip}{0pt}}
\usepackage{bookmark}
\IfFileExists{xurl.sty}{\usepackage{xurl}}{} % add URL line breaks if available
\urlstyle{same}
% fallback for those not using the hyperref driver hyperxmp:
\makeatletter
\@ifundefined{xmpquote}{\newcommand{\xmpquote}[1]{#1}}{}
\makeatother
\hypersetup{
  colorlinks=true,
  linkcolor={Maroon},
  filecolor={Maroon},
  citecolor={Blue},
  urlcolor={Blue},
  pdfcreator={LaTeX via pandoc}}

\author{}
\date{}

\begin{document}

\section{筑基篇课后习题
IX：物理空间的互锁结构------共享结构、丢失可观测、维度猜测}\label{ux7b51ux57faux7bc7ux8bfeux540eux4e60ux9898-ixux7269ux7406ux7a7aux95f4ux7684ux4e92ux9501ux7ed3ux6784ux5171ux4eabux7ed3ux6784ux4e22ux5931ux53efux89c2ux6d4bux7ef4ux5ea6ux731cux6d4b}

\begin{quote}
\textbf{声明 (Declaration)}: 本文\textbf{不代表任何数学结论}
(non-mathematical-claim), 仅为基于 Lean 4
形式化的一种\textbf{实验观测报告} (experimental observation report)。

{\def\LTcaptype{none} % do not increment counter
\begin{longtable}[]{@{}
  >{\raggedright\arraybackslash}p{(\linewidth - 6\tabcolsep) * \real{0.2500}}
  >{\raggedright\arraybackslash}p{(\linewidth - 6\tabcolsep) * \real{0.2500}}
  >{\raggedright\arraybackslash}p{(\linewidth - 6\tabcolsep) * \real{0.2500}}
  >{\raggedright\arraybackslash}p{(\linewidth - 6\tabcolsep) * \real{0.2500}}@{}}
\toprule\noalign{}
\begin{minipage}[b]{\linewidth}\raggedright
习题
\end{minipage} & \begin{minipage}[b]{\linewidth}\raggedright
主题
\end{minipage} & \begin{minipage}[b]{\linewidth}\raggedright
Lean 形式化
\end{minipage} & \begin{minipage}[b]{\linewidth}\raggedright
DOI
\end{minipage} \\
\midrule\noalign{}
\endhead
\bottomrule\noalign{}
\endlastfoot
exercise\_i\_9 & 物理空间的互锁结构（R166） &
PatPhysicalSpaceStructure.lean & 10.5281/zenodo.21916863 \\
\end{longtable}
}
\end{quote}

\textbf{Exercise IX for the Foundation-Building Chapter: The Interlock
Structure of Physical Space}

\begin{quote}
习题定位：筑基篇课后习题系列第九题。回答三个问题：①
五大力是否共享某些结构？② 丢失的结构中是否发生可由物理空间观测的数据？③
物理空间由多少互锁结构组成（节点/边/全互锁/维度/收缩）？全部新定理只锚筑基篇定理（R136/R138/R143/R147/R161/R163/R164/R165），不用外部引理。
\end{quote}

\emph{2026-08-13 · Internal research exercise · Lean 4 / mathlib v4.32.2
· 7 new theorems PROVED no sorry · 算器神魂论筑基篇课后习题 IX }

\begin{center}\rule{0.5\linewidth}{0.5pt}\end{center}

\subsection{习题陈述}\label{ux4e60ux9898ux9648ux8ff0}

\begin{quote}
① 这些力学（五大力），是否共享某些结构？②
这些力学，是否在物理空间丢失的结构中，发生了后续可由物理空间观测的数据？③
根据发现的结论，猜测物理空间由多少互锁结构组成，多少节点、多少边、是否全互锁、是几维结构、可能收缩了哪些维度？
\end{quote}

\begin{center}\rule{0.5\linewidth}{0.5pt}\end{center}

\subsection{一、★五大力共享互锁对单元（问题①）}\label{ux4e00ux4e94ux5927ux529bux5171ux4eabux4e92ux9501ux5bf9ux5355ux5143ux95eeux9898ux2460}

\textbf{★核心：五大力共享同一个基本单元------互锁对 \{d,
-d\}：exp(iθ)·exp(-iθ) = 1。}

\subsubsection{论证}\label{ux8bbaux8bc1}

\begin{enumerate}
\def\labelenumi{\arabic{enumi}.}
\tightlist
\item
  \textbf{互锁对 = 力的基本单元}（R136
  ②③：方向必须成对声明；R138：相位差可加；R143：对称对还原到
  1）：exp(iθ)·exp(-iθ) = 1（R161 pair\_interlock\_self\_consistent）。
\item
  \textbf{五大力的共享结构}：

  \begin{itemize}
  \tightlist
  \item
    力学：作用-反作用成对（F + (-F) = 0，R164 action\_reaction\_pair）
  \item
    量子：4 互锁 = S³ = SU(2)（R149/R154）
  \item
    电磁：E⊥B 正交（R047）
  \item
    强弱：规范群 k 对（R161）
  \item
    引力：脱离对逐对独立（R161）
  \end{itemize}
\item
  \textbf{共享单元}：每个力的动力学结构单元都是互锁对------成对的方向
  \{d, -d\}，组合还原（R143）。力 = 互锁对的破缺/组合/投影。
\end{enumerate}

\subsubsection{形式化（PatPhysicalSpaceStructure.lean，1
定理）}\label{ux5f62ux5f0fux5316patphysicalspacestructure.lean1-ux5b9aux7406}

\begin{itemize}
\tightlist
\item
  \texttt{forces\_share\_interlock\_pair}：exp(iθ)·exp(-iθ) =
  1------五大力共享互锁对单元。
\end{itemize}

\textbf{验证}：0 sorry。锚定 R161 pair\_interlock\_self\_consistent。

\begin{center}\rule{0.5\linewidth}{0.5pt}\end{center}

\subsection{二、★丢失结构的可观测性（问题②）}\label{ux4e8cux4e22ux5931ux7ed3ux6784ux7684ux53efux89c2ux6d4bux6027ux95eeux9898ux2461}

\textbf{★核心：一对互锁脱离物理空间后，剩余对仍互锁（物理法则保持）；脱离对的模长
‖ρ(g)‖ = 1 是脱离后仍可观测的不变量。}

\subsubsection{论证}\label{ux8bbaux8bc1-1}

\begin{enumerate}
\def\labelenumi{\arabic{enumi}.}
\tightlist
\item
  \textbf{脱离保持互锁}（R161
  pair\_detachment\_general：逐对独立）：一对互锁脱离物理空间（高维方向），剩余对仍互锁------\textbf{物理空间丢失结构后，剩余部分遵守物理空间法则}。
\item
  \textbf{单位圆模长 = 脱离不变观测}（R165
  phaseRep\_on\_circle/R141）：‖exp(iθ)‖ = 1
  对任意相位成立------脱离物理空间的对，其表示模长不随脱离改变，是\textbf{脱离后仍可观测的数据}（丢失结构的可观测印记）。
\item
  \textbf{丢失 → 可观测的机制}：脱离对进入高维（不可直接观测其相位
  θ），但其单位圆模长（=1）与剩余对的互锁不变量仍可观测------物理空间通过剩余结构''感知''丢失结构的存在（如引力
  = 脱离对的耦合效应，R164 CONJECTURE）。
\end{enumerate}

\subsubsection{形式化（PatPhysicalSpaceStructure.lean，2
定理）}\label{ux5f62ux5f0fux5316patphysicalspacestructure.lean2-ux5b9aux7406}

\begin{itemize}
\tightlist
\item
  \texttt{detachment\_keeps\_observable}：脱离某些对，剩余对仍互锁（逐对独立）。
\item
  \texttt{unit\_circle\_observable}：‖exp(iθ)‖ =
  1------脱离不变的单位圆模长观测。
\end{itemize}

\textbf{验证}：0 sorry。锚定 R161 / R165 phaseRep\_on\_circle。

\begin{center}\rule{0.5\linewidth}{0.5pt}\end{center}

\subsection{三、★物理空间结构猜测（问题③，CONJECTURE）}\label{ux4e09ux7269ux7406ux7a7aux95f4ux7ed3ux6784ux731cux6d4bux95eeux9898ux2462conjecture}

\textbf{★核心：物理空间 = 3 对共享互锁 = 3 节点 K3 完全图（3 边全互锁）=
3 维；可能收缩时间（R147 对合）与引力高维方向（R164 脱离对）。}

\subsubsection{论证}\label{ux8bbaux8bc1-2}

\begin{enumerate}
\def\labelenumi{\arabic{enumi}.}
\tightlist
\item
  \textbf{节点 = 3，边 = 3，全互锁}（R163：三体 = 3 节点全连接 = 三角形
  = K3 完全图，每对共享 2 相位 = 6 互锁）------物理空间的 3
  个独立方向两两全互锁。
\item
  \textbf{维度 = 3}（R161：3 对独立互锁 = 3 个方向自由度 = 6 相位 = ±x,
  ±y, ±z）------物理空间是 3 维。
\item
  \textbf{可能收缩的维度}：

  \begin{itemize}
  \tightlist
  \item
    \textbf{时间}（R147：时间 =
    对合对称对，未来/过去成对）------时间对可收缩（时间不是空间维度，但可观测其''相位''→
    因果/熵）
  \item
    \textbf{引力高维方向}（R164：引力 = 脱离投影，脱离对 =
    高维方向）------引力对应的对已脱离物理空间
  \item
    \textbf{完整结构可能 4 对}（3 空间 + 1 时间/引力），物理空间收缩到 3
    对（3 维）。
  \end{itemize}
\item
  \textbf{诚实边界}：这是结构猜测（CONJECTURE），依据
  R161/R163/R164/R147------``物理空间 = 3 对 K3 全互锁 3
  维''是观测归纳，非证明。
\end{enumerate}

\subsubsection{形式化（PatPhysicalSpaceStructure.lean，3
定理）}\label{ux5f62ux5f0fux5316patphysicalspacestructure.lean3-ux5b9aux7406}

\begin{itemize}
\tightlist
\item
  \texttt{physical\_space\_three\_pairs}：3 对独立互锁自洽（3
  节点，CONJECTURE 依据）。
\item
  \texttt{physical\_space\_K3\_complete}：K3 完全图（3 节点 3
  边全互锁，R163 三体结构）。
\item
  \texttt{physical\_space\_structure\_perspective}：★全景------3
  对共享互锁 ∧ K3 全互锁。
\end{itemize}

\textbf{验证}：0 sorry。锚定 R161 three\_independent\_pairs\_interlock /
R163 three\_body\_six\_shared\_interlock。

\begin{center}\rule{0.5\linewidth}{0.5pt}\end{center}

\subsection{四、习题解答总结}\label{ux56dbux4e60ux9898ux89e3ux7b54ux603bux7ed3}

{\def\LTcaptype{none} % do not increment counter
\begin{longtable}[]{@{}
  >{\raggedright\arraybackslash}p{(\linewidth - 6\tabcolsep) * \real{0.2500}}
  >{\raggedright\arraybackslash}p{(\linewidth - 6\tabcolsep) * \real{0.2500}}
  >{\raggedright\arraybackslash}p{(\linewidth - 6\tabcolsep) * \real{0.2500}}
  >{\raggedright\arraybackslash}p{(\linewidth - 6\tabcolsep) * \real{0.2500}}@{}}
\toprule\noalign{}
\begin{minipage}[b]{\linewidth}\raggedright
问题
\end{minipage} & \begin{minipage}[b]{\linewidth}\raggedright
回答
\end{minipage} & \begin{minipage}[b]{\linewidth}\raggedright
锚定
\end{minipage} & \begin{minipage}[b]{\linewidth}\raggedright
标签
\end{minipage} \\
\midrule\noalign{}
\endhead
\bottomrule\noalign{}
\endlastfoot
① 共享结构 & 五大力共享互锁对 \{d,-d\}（exp(iθ)·exp(-iθ)=1） &
R136②③/R138/R143/R161 & PROVED \\
② 丢失可观测 & 脱离保持互锁；单位圆模长 ‖ρ(g)‖=1 脱离不变 &
R161/R165/R141 & PROVED \\
③ 结构猜测 & 3 节点 3 边 K3 全互锁 3 维；收缩时间/引力对 &
R163/R161/R147/R164 & CONJECTURE \\
\end{longtable}
}

\textbf{习题的回答（三句话）}： 1.
\textbf{五大力共享互锁对单元}------每个力的基本结构都是成对的 \{d,
-d\}（exp(iθ)·exp(-iθ) = 1），力 = 互锁对的破缺/组合/投影。 2.
\textbf{丢失结构的可观测性}------脱离对进入高维后，其单位圆模长（=1）与剩余对的互锁不变量仍可观测；物理空间通过剩余结构''感知''丢失结构（引力
= 脱离对耦合效应）。 3. \textbf{物理空间 = 3 对共享互锁 = 3 节点 K3
完全图（3 边全互锁）= 3 维}------可能收缩了时间（R147
对合）与引力高维方向（R164 脱离对）；完整结构可能 4 对（3 空间 + 1
时间/引力）。

\begin{center}\rule{0.5\linewidth}{0.5pt}\end{center}

\subsection{五、相关对照}\label{ux4e94ux76f8ux5173ux5bf9ux7167}

{\def\LTcaptype{none} % do not increment counter
\begin{longtable}[]{@{}
  >{\raggedright\arraybackslash}p{(\linewidth - 4\tabcolsep) * \real{0.3333}}
  >{\raggedright\arraybackslash}p{(\linewidth - 4\tabcolsep) * \real{0.3333}}
  >{\raggedright\arraybackslash}p{(\linewidth - 4\tabcolsep) * \real{0.3333}}@{}}
\toprule\noalign{}
\begin{minipage}[b]{\linewidth}\raggedright
文献
\end{minipage} & \begin{minipage}[b]{\linewidth}\raggedright
对应内容
\end{minipage} & \begin{minipage}[b]{\linewidth}\raggedright
备注
\end{minipage} \\
\midrule\noalign{}
\endhead
\bottomrule\noalign{}
\endlastfoot
\textbf{Newton/Pauli/Maxwell/Yang-Mills}（习题 VII 对照） & 五大力 &
共享互锁对单元 \\
\textbf{R147}（筑基篇） & 时间 = 对合对称对 & 收缩维度的依据 \\
\textbf{R161/R163/R164/R165}（筑基篇） & k 对互锁/三体
K3/脱离投影/相位表示 & 结构猜测依据 \\
\end{longtable}
}

\begin{quote}
诚实边界：问题①（PROVED）与②（PROVED）为形式化结论；问题③（CONJECTURE）为结构猜测（3
节点 K3 3 维 + 收缩时间/引力对），依据框架结构归纳，非物理证明。
\end{quote}

\begin{center}\rule{0.5\linewidth}{0.5pt}\end{center}

\emph{筑基篇课后习题 IX · 2026-08-13 · Lean 4 / mathlib v4.32.2 · 7
theorems PROVED no sorry（PatPhysicalSpaceStructure.lean）· 配套 Lean
形式化：formal/Formal/Toolkit/PatPhysicalSpaceStructure.lean（快照见
../lean/PatPhysicalSpaceStructure.lean）}

\end{document}
